\chapter{Digital Modes and Modern Radio}
\label{ch:digital-modes}

Amateur radio has always adopted new technology quickly, and the past
two decades have brought a wave of digital modes and software-defined
techniques that meaningfully extend what a modest station can achieve.
This chapter surveys the modern digital operating landscape, building on
the modulation (Chapter~\ref{ch:modulation}) and error-control
(Chapter~\ref{ch:error-control}) foundations already covered.

\section{Why Digital Modes}

Compared to analog voice, digital modes can be engineered for
specific, quantifiable goals: minimum bandwidth for a given data rate,
maximum robustness to noise and fading at a given signal level, or
maximum throughput for file/data transfer. Because a digital signal's
information content is a defined, known set of discrete states, its
processing gain (Chapter~\ref{ch:error-control}) can be pushed far
beyond what human hearing achieves decoding analog voice or even CW by
ear, letting well-designed digital modes complete contacts at signal
levels where a voice or CW signal would be completely inaudible.

\section{Keyboard-to-Keyboard Text Modes}

\begin{itemize}
\item \textbf{RTTY (radioteletype)}: one of the oldest amateur digital
modes, using frequency-shift keying (Chapter~\ref{ch:modulation})
between two audio tones to send Baudot-coded text, historically
implemented with electromechanical teleprinters and now almost
universally via a computer sound card.
\item \textbf{PSK31}: a narrowband phase-shift-keyed mode designed
specifically for real-time keyboard-to-keyboard conversation, using a
variable-length coding scheme (shorter codes for more common letters)
that keeps typical typing speeds within an extremely narrow bandwidth
(around 31\,Hz, hence the name) --- allowing many simultaneous
PSK31 conversations within the bandwidth a single SSB voice signal
would occupy.
\end{itemize}

\section{Weak-Signal and Low-Rate Modes}

A newer generation of modes, largely built around combining strong
forward error correction (Chapter~\ref{ch:error-control}) with narrow
bandwidth and computer-controlled precise timing, trades data rate for
extreme sensitivity:

\begin{itemize}
\item \textbf{FT8} and related modes (part of the WSJT-X/MSHV software
family): structured, computer-synchronized exchanges (typically 15
seconds per transmission, tightly time-aligned to UTC) carrying a
minimal, highly compressed message set (callsigns, grid locators,
signal reports) using strong FEC coding, able to complete a logged
contact with signals well below the audible noise floor. The
constrained message format is a deliberate trade-off: extreme
sensitivity and speed of exchange, at the cost of the free-form
conversation possible on voice, CW, or PSK31.
\item \textbf{WSPR (Weak Signal Propagation Reporter)}: a very
low-rate, beacon-style mode designed not for two-way conversation but
for automatically probing and mapping propagation conditions, with
receiving stations automatically reporting decoded signals (including
estimated signal strength) to a shared online database.
\item \textbf{JT65 and related legacy modes}: earlier members of the
same weak-signal design family as FT8, historically important for
enabling reliable Earth-Moon-Earth (moonbounce) and other
extreme-path contacts that were previously impractical with analog
modes.
\end{itemize}

\section{Digital Voice}

Digital voice modes encode speech as compressed digital data (using a
\textbf{vocoder}) and transmit it much as any other digital data stream,
typically over FM. Several competing amateur digital voice systems
exist (differing in vocoder technology, network architecture, and
whether they are open or proprietary standards), generally offering
improved noise immunity and, where linked over the internet, wide-area
or global connectivity between repeaters and hotspots --- extending
local VHF/UHF coverage far beyond a single repeater's radio footprint,
though at the cost of a hard failure mode (digital signals tend to
either decode cleanly or not at all, rather than degrading gracefully
like analog FM near the noise floor).

\section{Packet Radio and Data Networking}

\textbf{Packet radio}, built on the amateur \textbf{AX.25} protocol
(itself derived from a wired data-networking standard, adapted for
radio), packages data into discrete, individually addressed and
checksummed packets (Chapter~\ref{ch:error-control}), enabling
store-and-forward messaging, position/status broadcasting (as used by
\textbf{APRS}, the Automatic Packet Reporting System, which
continuously broadcasts station location, weather, and status data over
a shared network of digipeaters and internet gateways), and
email-like messaging over radio-only paths (as used by systems such as
Winlink, particularly valuable for emergency and remote/maritime
communication where internet access is unavailable).

\section{Software-Defined Radio in Practice}

Building on the SDR receiver concepts introduced in
Chapter~\ref{ch:receivers}, modern amateur transceivers increasingly
implement transmit signal generation the same way: digital signal
processing generates the desired waveform (SSB, FM, digital mode audio,
or a fully digital signal) which is then converted to analog and
upconverted to the operating frequency, rather than being built up
through a chain of discrete analog mixer and filter stages. This
architecture makes it practical for a single piece of hardware to
support an enormous and continually growing range of modes purely
through software updates, and has made computer-controlled digital
modes far more accessible than in the era of dedicated
hardware modems for each individual mode.

\section{Choosing a Mode}

\begin{keyidea}
No single mode is universally ``best''; the right choice depends on the
goal. Voice (SSB/FM) suits real-time conversation and is the most
accessible to newcomers. CW remains highly effective in poor conditions
with minimal bandwidth and equipment. Weak-signal digital modes (FT8
and similar) excel at making contacts under marginal propagation but
support only minimal, structured exchanges. Packet-based data modes
(APRS, Winlink) suit position reporting and message delivery rather
than conversation at all. A well-rounded operator learns enough about
each family to select the right tool for a given goal and set of band
conditions.
\end{keyidea}

\begin{quickreview}
\item Digital modes trade among bandwidth, data rate, and sensitivity
in ways that can be engineered far more precisely than analog voice or
manual CW.
\item RTTY and PSK31 support real-time keyboard-to-keyboard text
conversation with modest bandwidth.
\item FT8, WSPR, and related weak-signal modes use strong FEC, narrow
bandwidth, and precise timing to decode signals below the audible noise
floor, at the cost of a constrained message format.
\item Digital voice modes offer noise-immune audio and, when linked,
wide-area connectivity, with a more abrupt (rather than gradual)
failure mode than analog FM.
\item Packet radio (AX.25), APRS, and Winlink extend amateur radio into
automated data networking, position reporting, and message delivery.
\end{quickreview}

\begin{practiceproblems}
\item Explain why FT8 can complete a contact when a voice signal at the
same power and antenna would be inaudible.
\item Explain the practical trade-off between FT8's extreme sensitivity
and its very limited message format.
\item Explain the purpose of APRS, and name one practical use beyond
simple vehicle tracking.
\item Explain why digital voice modes are described as having a more
``abrupt'' failure mode than analog FM as signal strength decreases.
\item Given a goal of making a contact under very poor HF propagation
conditions with a modest antenna, which mode family from this chapter
would you choose, and why?
\end{practiceproblems}
